\section{Step 95 Write both outputs that cross input boundaries}
Fill \texttt{pack\_shuffle\_middle}. Produce output words four through seven using input vectors zero and one, then words eight through eleven using vectors one and two. Restore \texttt{\$w3} from the appropriate retained mask before each shuffle, and store at destination offsets sixteen and thirty-two.

For the first crossing, place input zero on the right and input one on the left, as in the earlier worked shuffle. Selectors zero through fifteen address the right operand; sixteen through thirty-one address the left. For the second crossing, the mask’s origin moves forward by sixteen global bytes because the operand pair is now one and two.

At \texttt{pack\_middle\_written}, the first twelve output words match the public RGB samples, while words twelve through fifteen remain untouched. The mask table and retained vectors still have their original values. Swapping the two input operands can produce plausible colors in some lanes while corrupting precisely these boundary pixels.

\paragraph{Observe} Use \texttt{s95-1}; stop before \texttt{pack\_middle\_written}. Inspect the state described above before continuing.

\Needspace{7\baselineskip}
\paragraph{Check} For the second boundary shuffle, what is the mask’s global-byte origin?
\begin{itemize}
\item[A.] Zero.
\item[B.] Thirty-two.
\item[C.] Sixteen.
\end{itemize}
